\section{Primitive factors and bounded classification}\label{sec:bounded}

\subsection{Cycle and theta cores}

Move a root-containing factor along an all-tree path to a real labelled
incoming boundary, using \cref{lem:root-movement}.  Delete port arms and
suppress ordinary bivalent vertices that are neither the local source nor a
reticulation sink, recording the removed port-bearing subdivisions as
ordered words.

\begin{proposition}[Cycle--theta primitive theorem]\label{prop:cores}
Every nontrivial complete blob factor in the strong level-2 class has a cycle
core or one of four directed theta cores
\(\theta_0,\theta_1,\theta_2,\theta_3\).  It is recovered uniquely from the
core, one labelled ordered port word on each directed segment, and one child
port below every path-sink reticulation.
\end{proposition}

\begin{proof}
Let \(r,v,e\) be the numbers of reticulations, internal vertices, and blob
edges after port arms are removed.  Summing indegrees gives
\(e=v+r-1\), hence

\[
 \sum_w(\deg_B(w)-2)=2e-2v=2(r-1).
\]

For \(r=1\), every blob vertex has degree two and the core is a cycle.  For
\(r=2\), exactly two vertices have degree three; biconnectedness makes the
blob three internally disjoint paths between them, hence a theta.  The
incoming boundary is ordinary: deleting an arc entering a reticulation cannot
disconnect its endpoints, because the other parent and their root paths give
an alternate connection.  Its factor endpoint is therefore the unique local
tree source \(S\).  A theta pole already uses all three binary incidences in
the blob, so \(S\) lies inside one pole-to-pole path.  Any internal
reticulation is a path sink whose unique child is its boundary arm.

At most one pole is reticulate.  If both were, the path containing \(S\)
points toward both poles and cannot contain either pole's unique outgoing
side.  Those two outgoing sides must therefore lie on the other event-free
paths, forcing one complete path from the first pole to the second and the
other back from the second to the first.  That is a directed cycle.

Suppose exactly one pole \(R\) is reticulate and the other \(T\) is a tree
pole.  The second reticulation is an internal sink \(X\).  If \(S\) and \(X\)
lie on different paths, the \(S\)-path points away from \(S\), the
\(X\)-path points toward \(X\), and the third path is forced from \(T\) to
\(R\); this is \(\theta_0\).  If they share a path, the only acyclic order
from \(T\) to \(R\) is \(T\)--\(S\)--\(X\)--\(R\).  The reverse event order
would force the two event-free paths in opposite directions and create a
directed cycle.  Both event-free paths therefore point from \(T\) to \(R\),
giving \(\theta_1\).

If both poles are tree vertices, both reticulations are internal sinks.  They
cannot share one path: the two arrows entering each sink would force \(S\)
between the sinks, while each of the other two event-free pole-to-pole paths
would then have to be directed from one tree pole to the other.  Those paths
cannot have the same direction, because the receiving pole would have two
parents; opposite directions form a directed cycle with the
source-to-sink subpaths.  Thus the sinks lie on distinct paths.  If \(S\)
lies on the third path, both poles feed both sinks, giving \(\theta_2\).
Otherwise \(S\) shares a path with one sink; read from the pole whose other
two sides point outward, acyclicity fixes the order
pole--\(S\)--sink--pole and every remaining direction, giving
\(\theta_3\).  Reversing unnamed poles or interchanging unnamed sinks yields
only the displayed symmetries.  Hence no source or sink placement is omitted.
This is the directed ported refinement of the two underlying
level-2 generators in \cite{HuberEtAl2025}; the frozen topology replay is a
finite regression of this case split, not its premise.
\end{proof}

\begin{figure}[t]
\centering
\begin{tikzpicture}[
  x=1cm,y=1cm,
  every node/.style={font=\scriptsize},
  source/.style={circle,draw,inner sep=1pt,minimum size=4.2mm},
  tree/.style={circle,draw,inner sep=1pt,minimum size=4.2mm},
  retic/.style={circle,draw,double,inner sep=.85pt,minimum size=4.7mm},
  seg/.style={-{Stealth[length=1.55mm,width=1.05mm]},line width=.52pt},
  port/.style={seg,dashed},
  slabel/.style={fill=white,inner sep=.55pt,font=\tiny},
  title/.style={font=\small},
  repair/.style={font=\tiny,align=center}
]
  % Cycle.
  \begin{scope}[shift={(0,4.0)}]
    \node[title] at (-1.0,1.25) {cycle};
    \node[source] (s) at (0,.75) {$S$};
    \node[retic] (x) at (0,-.75) {$X$};
    \draw[port] (0,1.35) -- node[right,slabel] {$\iota$} (s);
    \draw[seg] (s) to[bend right=42] node[left,slabel] {$0$} (x);
    \draw[seg] (s) to[bend left=42] node[right,slabel] {$1$} (x);
    \draw[port] (x) -- (0,-1.35);
    \node[repair] at (0,-1.72) {$\mathcal R_{\min}:\{0\},\{1\}$};
  \end{scope}

  % theta_0.
  \begin{scope}[shift={(4.25,4.0)}]
    \node[title] at (-1.45,1.25) {$\theta_0$};
    \node[source] (s) at (0,.8) {$S$};
    \node[tree] (u) at (-1,0) {$U$};
    \node[retic] (v) at (1,0) {$V$};
    \node[retic] (x) at (0,-1) {$X$};
    \draw[port] (0,1.38) -- node[right,slabel] {$\iota$} (s);
    \draw[seg] (s) -- node[left,slabel] {$0$} (u);
    \draw[seg] (s) -- node[right,slabel] {$1$} (v);
    \draw[seg] (u) -- node[left,slabel] {$2$} (x);
    \draw[seg] (v) -- node[right,slabel] {$3$} (x);
    \draw[seg] (u) -- node[above,slabel] {$4$} (v);
    \draw[port] (x) -- (0,-1.48);
    \node[repair] at (0,-1.78)
      {$\mathcal R_{\min}:\{2,3\},\{3,4\}$};
  \end{scope}

  % theta_1.
  \begin{scope}[shift={(9.0,4.0)}]
    \node[title] at (-1.55,1.25) {$\theta_1$};
    \node[source] (s) at (-1,.75) {$S$};
    \node[retic] (x) at (1,.75) {$X$};
    \node[tree] (u) at (-1,-.75) {$U$};
    \node[retic] (v) at (1,-.75) {$V$};
    \draw[port] (-1,1.35) -- node[right,slabel] {$\iota$} (s);
    \draw[seg] (s) -- node[left,slabel] {$0$} (u);
    \draw[seg] (s) -- node[above,slabel] {$1$} (x);
    \draw[seg] (v) -- node[right,slabel] {$2$} (x);
    \draw[seg] (u) to[bend left=24] node[above,slabel] {$3$} (v);
    \draw[seg] (u) to[bend right=24] node[below,slabel] {$4$} (v);
    \draw[port] (x) -- (1.65,.75);
    \node[repair] at (0,-1.52)
      {$\mathcal R_{\min}:\{2,3\},\{2,4\}$};
  \end{scope}

  % theta_2.
  \begin{scope}[shift={(2.2,.25)}]
    \node[title] at (-1.7,1.35) {$\theta_2$};
    \node[source] (s) at (0,1.0) {$S$};
    \node[tree] (u) at (-1.15,.15) {$U$};
    \node[tree] (v) at (1.15,.15) {$V$};
    \node[retic] (x0) at (-.52,-1.0) {$X_0$};
    \node[retic] (x1) at (.52,-1.0) {$X_1$};
    \draw[port] (0,1.58) -- node[right,slabel] {$\iota$} (s);
    \draw[seg] (s) -- node[left,slabel] {$0$} (u);
    \draw[seg] (s) -- node[right,slabel] {$1$} (v);
    \draw[seg] (u) -- node[left,slabel] {$2$} (x0);
    \draw[seg] (v) -- node[pos=.34,right,slabel] {$3$} (x0);
    \draw[seg,preaction={draw=white,line width=2pt,-}]
      (u) -- node[pos=.34,left,slabel] {$4$} (x1);
    \draw[seg] (v) -- node[right,slabel] {$5$} (x1);
    \draw[port] (x0) -- (-.52,-1.48);
    \draw[port] (x1) -- (.52,-1.48);
    \node[repair] at (0,-1.83)
      {$\mathcal R_{\min}:\{2,3\},\{2,5\},\{3,4\},\{4,5\}$};
  \end{scope}

  % theta_3.
  \begin{scope}[shift={(7.4,.25)}]
    \node[title] at (-1.7,1.35) {$\theta_3$};
    \node[source] (s) at (0,1.0) {$S$};
    \node[tree] (u) at (-1.25,.15) {$U$};
    \node[retic] (x0) at (1.25,.15) {$X_0$};
    \node[tree] (v) at (0,-.65) {$V$};
    \node[retic] (x1) at (-1.25,-1.35) {$X_1$};
    \draw[port] (0,1.58) -- node[right,slabel] {$\iota$} (s);
    \draw[seg] (s) -- node[left,slabel] {$0$} (u);
    \draw[seg] (s) -- node[right,slabel] {$1$} (x0);
    \draw[seg] (v) -- node[right,slabel] {$2$} (x0);
    \draw[seg] (u) -- node[left,slabel] {$3$} (x1);
    \draw[seg] (v) -- node[below,slabel] {$4$} (x1);
    \draw[seg] (u) -- node[above,slabel] {$5$} (v);
    \draw[port] (x0) -- (1.75,.15);
    \draw[port] (x1) -- (-1.25,-1.72);
    \node[repair] at (.25,-1.98)
      {$\mathcal R_{\min}:\{2\},\{4\}$};
  \end{scope}

\end{tikzpicture}
\caption{The five primitive complete-factor cores.  Segment indices agree
with the port-word grammar; $\iota$ is the selected incoming boundary.
The complete minimum strong repairs $\mathcal R_{\min}$ are shown below
each core.  Single circles mark source or tree roles, double circles mark
reticulations, and dashed stubs retain incoming or path-sink child ports.}
\label{fig:primitive-cores}
\end{figure}


Figure~\ref{fig:primitive-cores} fixes the pole, segment, incoming-boundary,
and minimum-repair notation used in the finite grammar.

The directed segments and minimum occupied sets required by
\cref{lem:no-omnian} are:

\begin{center}
\small
\begin{tabular}{lll}
\toprule
core & directed segments in index order & minimum occupied sets\\
\midrule
cycle & \(S\to X,S\to X\) & \(\{0\},\{1\}\)\\
\(\theta_0\) & \(S\to U,S\to V,U\to X,V\to X,U\to V\)
 & \(\{2,3\},\{3,4\}\)\\
\(\theta_1\) & \(S\to U,S\to X,V\to X,U\to V,U\to V\)
 & \(\{2,3\},\{2,4\}\)\\
\(\theta_2\) & \(S\to U,S\to V,U\to X_0,V\to X_0,\)
 \(U\to X_1,V\to X_1\)
 & \(\{2,3\},\{2,5\},\{3,4\},\{4,5\}\)\\
\(\theta_3\) & \(S\to U,S\to X_0,V\to X_0,U\to X_1,\)
 \(V\to X_1,U\to V\)
 & \(\{2\},\{4\}\)\\
\bottomrule
\end{tabular}
\end{center}

The corresponding obstruction clauses are
\(\{0,1\}\) for a cycle; \(\{3\},\{2,4\}\) for \(\theta_0\);
\(\{2\},\{3,4\}\) for \(\theta_1\);
\(\{2,4\},\{3,5\}\) for \(\theta_2\); and \(\{2,4\}\) for
\(\theta_3\).  The displayed sets are precisely their minimum transversals.

For a target core with \(m_H\) directed segments, \(q_H\) path-sink roles,
and \(r_H\) repair tags, the exact number of target completions using \(k\)
selected boundaries and selected-incoming flag \(\varepsilon\) is

\begin{equation}\label{eq:completion-count}
 C(k,\varepsilon)=\sum_H r_H\sum_{j=0}^{q_H}
 \binom{q_H}{j}
 \binom{k-\varepsilon-j+m_H-1}{m_H-1}.
\end{equation}

The tuples \((m_H,q_H,r_H)\) are
\((2,1,1),(5,1,2),(5,1,2),(6,2,4),(6,2,2)\).  Thus the four-port target
grammar has 831 selected-incoming and 1,983 marginalized-incoming
completions.

\subsection{The exact four-port residue}

Six rigid K3P source supports have normalized ranks

\[
 20,\ 21,\ 21,\ 21,\ 23,\ 24.
\]

After pointwise, quadratic, graph-isomorphism, and ordinary-triangle filters,
the only nonterminal four-port residue consists of the following 14
canonical directed orbits.  Here the arrow in the rank column is source to
target, and ``H14'' means a transported marginal of \eqref{eq:H14}.  The
symbol \(r_j\) is the zero-based index of the listed minimum occupied set for
that core, and a word \texttt{p0p1p2p3} records the port relabeling
\(i\mapsto p_i\).  The family labels H21 and L\(d\) record, respectively, a
rank-21 nonautomorphic relabeling family and a lower-rank source mapping to a
rank-24 target; the suffixes a and b distinguish the two
\(\theta_1\)-repair families.

{\small
\begin{longtable}{llllrl}
\toprule
orbit & source & target & rank & raw & exact certificate\\
\midrule
\endhead
H21-01 & \(\theta_0,r_1\) & \(\theta_0,r_1;\texttt{0132}\) & \(21\to21\) & 4 & H14 quartic\\
H21-02 & \(\theta_0,r_1\) & \(\theta_0,r_1;\texttt{0213}\) & \(21\to21\) & 4 & \(11>10\) rank\\
H21-03 & \(\theta_0,r_1\) & \(\theta_0,r_1;\texttt{0231}\) & \(21\to21\) & 4 & H14 quartic\\
H21-04 & \(\theta_0,r_1\) & \(\theta_0,r_1;\texttt{0312}\) & \(21\to21\) & 4 & H14 quartic\\
H21-05 & \(\theta_0,r_1\) & \(\theta_0,r_1;\texttt{0321}\) & \(21\to21\) & 2 & H14 quartic\\
H21-06 & \(\theta_0,r_1\) & \(\theta_0,r_1;\texttt{1032}\) & \(21\to21\) & 4 & H14 quartic\\
L20-01 & \(\theta_0,r_0\) & \(\theta_3,r_1;\texttt{1203}\) & \(20\to24\) & 2 & quartic\\
L20-02 & \(\theta_0,r_0\) & \(\theta_3,r_1;\texttt{1230}\) & \(20\to24\) & 2 & \(14>12\) rank\\
L21a-01 & \(\theta_1,r_0\) & \(\theta_3,r_1;\texttt{0123}\) & \(21\to24\) & 2 & quartic\\
L21a-02 & \(\theta_1,r_0\) & \(\theta_3,r_1;\texttt{0132}\) & \(21\to24\) & 2 & \(11>10\) rank\\
L21b-01 & \(\theta_1,r_1\) & \(\theta_3,r_1;\texttt{0123}\) & \(21\to24\) & 2 & quartic\\
L21b-02 & \(\theta_1,r_1\) & \(\theta_3,r_1;\texttt{0132}\) & \(21\to24\) & 2 & \(11>10\) rank\\
L23-01 & \(\theta_3,r_0\) & \(\theta_3,r_1;\texttt{0123}\) & \(23\to24\) & 2 & \(14>12\) rank\\
L23-02 & \(\theta_3,r_0\) & \(\theta_3,r_1;\texttt{0132}\) & \(23\to24\) & 2 & quartic\\
\bottomrule
\end{longtable}
}

The two pre-lock sink swaps are separate from those fourteen orbits.  Their
common rigid support is \(\theta_3,r_1\) with selected incoming port, both
sink-child ports, and the repair port on segment 4; the rows below are its
two nonautomorphic labelled presentations.

\begin{center}
\small
\begin{tabular}{llll}
\toprule
presentation & port relabeling & rank & exact certificate\\
\midrule
SINK-0132 & \texttt{0132} & \(24\to24\) & quartic\\
SINK-1023 & \texttt{1023} & \(24\to24\) & quartic\\
\bottomrule
\end{tabular}
\end{center}

The six H21 rows are the nonisomorphic double cosets of the repaired
\(\theta_0\) source; the seventh double coset is isomorphic.  The complete
final post-filter accounting is

\begin{equation}\label{eq:38plus2}
 \boxed{40=38+2.}
\end{equation}

The 38 raw records belong to the fourteen named orbits; the two separately
tabulated rows are rank-\(24\) \(\theta_3\) sink-swap presentations.  Neither is a
labelled graph isomorphism, and the same exact eight-term quartic has zero
target pullback and nonzero strict
source value on each.

\begin{theorem}[Fourteen-orbit classification]\label{thm:fourteen}
The polynomially separated orbits are

\[
\begin{gathered}
\text{H21-01, H21-03, H21-04, H21-05, H21-06,}\\
\text{L20-01, L21a-01, L21b-01, L23-02},
\end{gathered}
\]

and the directed-rank-separated orbits are

\[
 \text{H21-02, L20-02, L21a-02, L21b-02, L23-01}.
\]

The two sink swaps are polynomially separated.  Thus the bounded residue has
zero unresolved directions, no new symmetric K3P move, and no proper
directed K3P containment.
\end{theorem}

\begin{proof}[Exact finite proof]
The active exhaustive replay begins with the six source supports, the 2,814
target completions obtained from \eqref{eq:completion-count}, and all 24 port
permutations.  It visits all
\(6\cdot2814\cdot24=405{,}216\) directed presentations.  Literal graph
tests exclude 377,382, leaving 27,834 post-topology presentations.  Exact
Jacobian minors together with coefficientwise polynomial-vector-field upper
bounds exclude 23,054; the quadratic identities exclude 1,968; transported
\(H_{14}\) marginals exclude 88; and 2,540 enter the independently replayed
restoration forest.  The remaining presentations comprise 30 labelled
isomorphisms, 114 ordinary-triangle presentations, and the 40 records in
\eqref{eq:38plus2}.  Thus this accounting is derived from the primitive
grammar rather than read from the fourteen-orbit lock.

For a literal target map \(f\), let \(A\) be the exact coefficient matrix of
the polynomial vector-field identities \(J_fV=0\), and let \(E\) evaluate the
vector-field coefficients at the recorded strict rational point.  Exact linear
algebra computes

\[
 \dim E(\ker A)=\rank\!\begin{bmatrix}A\\E\end{bmatrix}-\rank(A).
\]

A nonzero evaluation-image minor makes these kernel fields generically
independent and hence gives the certified target-rank upper.  Sampled Jacobian
ranks are used only to supply nonzero lower-rank minors, never as upper bounds.
A presentation enters the 23,054-case rank-exclusion class only when its exact
source lower rank is strictly greater than that certified target upper.

The separately implemented verifier imports neither the producer nor its
graph-atlas module.  It reconstructs the grammar, all 4,379 distinct literal
maps used by the presentations, the topology tests, every exact nonzero rank
minor, every target rank upper bound, and the polynomial pullbacks.  It also
checks the 2,540 restoration obligations against the active forest by their
source, target, and port-permutation data, and derives the double-coset
quotient of the 40-record residue.  Coherently resealed omission and
reclassification mutations are required to fail.

For each of the nine polynomial orbits, an exact characteristic-zero quartic
vanishes on the complete target map and has a nonzero value at a strict
rational source point.  The first five are transported \(H_{14}\) marginal
quartics; the remaining four are multigraded eight-term quartics.  Exact graph
and Fourier transports cover every raw member.

For each of the five rank orbits, a selected source projection has a nonzero
exact Jacobian minor of rank \(11\) or \(14\).  The target projection factors
rationally through 10 or 12 generators, respectively, giving the
strict upper bound.  These are directional selected-projection obstructions,
not comparisons of total model dimension: the target may have equal or
larger total normalized rank while its indicated observable projection has
smaller function-field dimension.  The H21 upper bound holds after saturation
by declared physical factors, all nonzero throughout the strict domain; the
other four targets are compressed ordinary-sunlet projections.  A
source-relative open containment cannot lie inside the lower-dimensional
target projection.

The clean-room implementation reconstructs all six sources, all 2,814
target completions, seven H21 double cosets, all 38 coordinate transports,
and both sink swaps.  The original H21-01 failure is diagnosed as the use of
rooted-DAG automorphisms in place of root-suppressed semi-directed mixed-graph
automorphisms, compounded by failure to conjugate a base target automorphism
into the displayed frame.  The corrected mixed-graph transport passes and
the historical verifier remains preserved.  Exact rows, columns, bodies,
source evaluations, and hashes appear in the supplement.
\end{proof}

The bounded result is intentionally direction-sensitive.  A target of larger
generic dimension is not automatically a containing model, and a polynomial
identity is transported through the literal labelled graph rather than
attached only to a topology name.
